% Insert as A.40 after the full class-two lift A.38 and recognition A.39.
\begin{statement}{A.40}{Proved}{The twenty-four central twists and the exact local matching map}
Use the actual class-two quotient $E_2=H(\widehat\Z)\rtimes G_\Q$
and its coordinates from A.38.  There are isomorphisms
\[
 H^1_{\rm cts}(G_\Q,\Z_\ell(2))\simeq
 \begin{cases}
  \Z/8,&\ell=2,\\
  \Z/3,&\ell=3,\\
  0,&\ell\ge5,
 \end{cases}
 \qquad
 H^1_{\rm cts}(G_\Q,\widehat\Z(2))\simeq\Z/24.
 \tag{A.40.1}
\]
In the chosen Tate coordinates a generator of the last group is
the continuous cocycle
\[
                \eta(g)=\frac{\chi(g)^2-1}{24}.
 \tag{A.40.2}
\]
For $(A,B)\in\mathcal G$, the geometric conjugacy classes of
sections of $E_2\to G_\Q$ above its fixed abelian section class
form a torsor under this group.  Thus this fiber contains exactly
twenty-four classes.

The localization homomorphism on these central twists is
\[
 \operatorname{loc}:\Z/24\longrightarrow
       \prod_{q<\infty}H^1_{\rm cts}(G_{\Q_q},\widehat\Z(2)),
 \qquad k\longmapsto(\operatorname{res}_q(k\eta))_q.
 \tag{A.40.3}
\]
It is injective.  Its composite with the two coordinate projections
at $(q,\ell)=(2,3),(3,2)$ is the isomorphism
\[
 \begin{split}
 \Z/24&\ \xrightarrow{\ \sim\ }
       \Z_3/3\Z_3\oplus\Z_2/8\Z_2,\\
 k&\longmapsto(k/8\bmod3,\ k/3\bmod8).
 \end{split}
 \tag{A.40.4}
\]

For each ordinary local coordinate $x_q$ of $(A,B)$, the actual
local point section determines a local comparison class
$\delta_q$ relative to any chosen global class-two lift, as
constructed in the proof.  The subset of global lift classes
agreeing with all these local point-section classes is exactly
\[
 \bigl\{s+k\eta:\ k\in\Z/24,
              \ (\delta_q)_q=\operatorname{loc}(k)\bigr\}.
 \tag{A.40.5}
\]
Consequently this subset has at most one element.  For every
$(A,B)\in\mathcal G$, there is a unique global class-two lift
whose two central comparison coordinates in (A.40.4) agree with
the local points at $2$ and $3$.  Formula (A.40.5) is the exact
test for that distinguished lift to agree with the point sections
at all finite places.
\end{statement}
\begin{proof}
\emph{Descent of every integral one-cocycle to a fixed finite
ramification set.}
Fix $\ell$ and put $M=\Z_\ell(2)$ and $S_\ell=\{\ell,\infty\}$.
At a prime $q\ne\ell$, inertia acts trivially on $M$.  The
restriction of a continuous one-cocycle to inertia is therefore
a continuous homomorphism to the additive torsion-free pro-$\ell$
group $M$.  It kills wild inertia and all prime-to-$\ell$ tame
factors.  Write $t$ for its value on a generator of the remaining
tame $\Z_\ell(1)$, and $F$ for an arithmetic Frobenius.  The
relation $F\tau F^{-1}=\tau^q$ and the cocycle identity give
\[
                         q^2t=qt.
 \tag{A.40.6}
\]
Since $M$ is torsion-free, $t=0$.  The cocycle vanishes on every
such inertia group.  It also vanishes on their conjugates, because
these groups act trivially on $M$, and hence on their closed normal
closure.  Inflation therefore gives an isomorphism
\[
 H^1_{\rm cts}(G_{\Q,S_\ell},M)
       \xrightarrow{\ \sim\ }H^1_{\rm cts}(G_\Q,M).
 \tag{A.40.7}
\]
This argument applies to the cocycles themselves, not just their
finite-coefficient classes.

\emph{Finiteness, including the real correction at two.}
Set $G_S=G_{\Q,S_\ell}$ and $M_n=M/\ell^nM$.  The unramified
calculations in A.27 give
\[
 H^2_{\rm cts}(G_S,M)=
 \begin{cases}\Z/2,&\ell=2,\\0,&\ell\text{ odd}.\end{cases}
 \tag{A.40.8}
\]
Here the group at two is detected by restriction to the real
place.  We also need $H^3_{\rm cts}(G_S,M)=0$, and verify it
before using a coefficient exact sequence.  All finite-coefficient
cohomology groups are finite by
\cite[I, Corollary 4.15]{BIR_MilneADT}.  The comparison with
continuous cochains and its inverse-limit exact sequence is
\cite[\S1, (1.6), and Theorem 2.2]{BIR_Jannsen1988}.
The finiteness makes the relevant $\varprojlim^1$ terms zero.
By \cite[I, Theorem 4.10(c)]{BIR_MilneADT},
\[
                 H^3(G_S,M_n)=H^3(G_\R,M_n).
\]
Complex conjugation acts trivially on the even twist.  For odd
$\ell$ the last group is zero.  For $\ell=2$ it is
$M_n[2]=\{0,2^{n-1}\}$, and every reduction map from level
$n+1$ to level $n$ is zero on this group.  Its inverse limit is
therefore zero in this case as well.  This proves the required
integral $H^3$ vanishing.

The exact sequence $0\to M\xrightarrow{\ell^n}M\to M_n\to0$
is exact on continuous cochains: a continuous finite-valued
cochain has a locally constant lift to $M$.  Its cohomology
sequence and (A.40.8) now give
\[
 H^2(G_S,M_n)=
 \begin{cases}\Z/2,&\ell=2,\\0,&\ell\text{ odd}.\end{cases}
 \tag{A.40.9}
\]
The global Euler characteristic formula
\cite[I, Theorem 5.1]{BIR_MilneADT} has archimedean factor one
here: at the real place $H^0(G_\R,M_n)=M_n$, whose order equals
the normalized real absolute value of $\#M_n$.  In particular
the formula applies at two with its stated ordinary $H^0$ and
truncated Euler characteristic; no assumption that the real
higher cohomology vanishes is being made.  It gives
\[
          \#H^1(G_S,M_n)=\#H^0(G_S,M_n)\,\#H^2(G_S,M_n).
 \tag{A.40.10}
\]
The cyclotomic image of $G_S$ is $\Z_\ell^\times$.  The ideal
generated by its elements $u^2-1$ is
\[
 I_\ell=(u^2-1:\ u\in\Z_\ell^\times)=
 \begin{cases}
  8\Z_2,&\ell=2,\\
  3\Z_3,&\ell=3,\\
  \Z_\ell,&\ell\ge5.
 \end{cases}
 \tag{A.40.11}
\]
Indeed odd squares are $1$ modulo $8$ and $3^2-1=8$;
at three all nonzero squares are $1$ modulo $3$ and $2^2-1=3$;
at $\ell\ge5$ a nonzero residue different from $\pm1$ gives a
unit $u^2-1$.  Thus the orders of $H^0(G_S,M_n)$ are bounded by
$8,3,1$, respectively.  Equations (A.40.9)--(A.40.10) bound the
orders of the finite $H^1$ groups by $16,3,1$.  Since
\[
 H^1_{\rm cts}(G_S,M)=\varprojlim_nH^1(G_S,M_n)
\]
and these finite groups have uniformly bounded order, the
integral $H^1$ group is finite.

\emph{Its exact order and explicit generator.}
The preceding finiteness is sufficient to compute the group;
the bound $16$ at two is not its order.  If a cocycle $c$ has
class killed by $\ell^r$, then $\ell^r c(g)=(\chi(g)^2-1)t$
for some $t\in M$.  Consequently
\[
                   c(g)=(\chi(g)^2-1)(t/\ell^r)
\]
inside $\Q_\ell(2)$.  Integrality says precisely that the class
of $t/\ell^r$ in $(\Q_\ell/\Z_\ell)(2)$ is $G_S$-invariant.
Conversely every such invariant gives this integral cocycle.
Its class is zero only when the fractional element differs from
an element of $M$ by an invariant of $\Q_\ell(2)$; the latter
invariant group is zero.  Thus the connecting map gives the
explicit isomorphism
\[
 H^0(G_S,(\Q_\ell/\Z_\ell)(2))
    \xrightarrow{\ \sim\ }H^1_{\rm cts}(G_S,M),
 \quad \bar r\longmapsto[g\mapsto(\chi(g)^2-1)r].
 \tag{A.40.12}
\]
By (A.40.11), its source has order $8,3,1$ respectively.
This proves the prime-by-prime assertion in (A.40.1), including
the exact order eight at two.

As in (A.38.9), continuous cochains with values in
$\widehat\Z(2)=\prod_\ell\Z_\ell(2)$ are the product of the
coefficient cochain complexes, and taking cohomology commutes
with this product.  The full group is therefore
$\Z/8\oplus\Z/3$.  The numerator of (A.40.2) is divisible by
$24$ in $\widehat\Z$ by (A.40.11), so its displayed quotient
is integral and continuous.  It satisfies
\[
        \eta(gh)=\eta(g)+\chi(g)^2\eta(h).
\]
At two the element $1/24$ generates the invariant fractional
quotient of order eight; at three it generates that of order
three; at every other coefficient prime it is integral and
gives a coboundary.  Hence $[\eta]$ has order twenty-four and
gives the stated isomorphism.

\emph{Local coordinates and injectivity.}
The argument (A.40.6) also proves that every integral cocycle
on $G_{\Q_q}$ with values in $\Z_\ell(2)$ is unramified when
$q\ne\ell$.  For the procyclic unramified quotient, evaluation
at arithmetic Frobenius identifies cocycles modulo coboundaries
with
\[
        H^1_{\rm cts}(G_{\Q_q},\Z_\ell(2))
                =\Z_\ell/(q^2-1)\Z_\ell,\qquad q\ne\ell.
 \tag{A.40.13}
\]
Every value in $\Z_\ell$ occurs: modulo $\ell^n$, the cocycle
on integral powers of Frobenius is periodic after a multiple of
$\ell^n$ times the order of $q^2$ modulo $\ell^n$; these
compatible finite cocycles extend to the procyclic group.  A
coboundary has Frobenius value in $(q^2-1)\Z_\ell$.
Formula (A.40.2) therefore restricts to the class
\[
                       \frac{q^2-1}{24}
             \pmod{(q^2-1)\Z_\ell}.
 \tag{A.40.14}
\]
Taking $(q,\ell)=(2,3)$ and $(3,2)$ gives (A.40.4).
The two displayed elements generate the respective cyclic
groups, so the Chinese remainder theorem proves that the
composite is an isomorphism.  In particular (A.40.3) is
injective.

\emph{The actual fiber of the class-two section map.}
Choose continuous cocycles $a,b$ for $A,B$ and a lift from
A.38, written
\[
                   s(g)=(a(g),b(g),h(g);g).
\]
Every other section with exactly these abelian cocycles has
third coordinate $h+c$ for a continuous cocycle
$c:G_\Q\to\widehat\Z(2)$, since the two third-coordinate
defects are the same cup cocycle.  Conjugation by a central
element $(0,0,z)$ changes $c$ by the coboundary
$g\mapsto z-\chi(g)^2z$.  Thus these lifts modulo central
conjugation form a torsor under the group (A.40.1).

This is also the fiber on geometric conjugacy classes, with
no extra quotient implicit in the assertion.  In the abelian
quotient, conjugation by an element with geometric coordinates
$(u,v)$ changes $(a,b)$ by
$((1-\chi)u,(1-\chi)v)$.  Both globally and over every
$\Q_q$,
\[
                  H^0(G,\widehat\Z(1)^2)=0.
 \tag{A.40.15}
\]
For $\ell\ne q$ the local assertion follows from the nonzero
scalar $q-1$ acting on the torsion-free $\Z_\ell$ factor;
for $\ell=q$ the cyclotomic image contains a unit other than
one, and the same conclusion holds.  The global assertion
follows from the surjectivity of each cyclotomic character.
Consequently a geometric conjugation preserving the chosen
abelian cocycles has $u=v=0$ and is central.  Any section in
the fiber can first be conjugated to have the chosen abelian
cocycles.  The required abelian conjugating element is unique
by (A.40.15), and any two of its lifts differ centrally.
This proves the torsor assertion on the stated conjugacy classes.

For a local point section at $x_q$, its abelian cohomology
classes are the restrictions of $A,B$.  Choose a geometric
conjugation making its actual abelian cocycles equal to
$a|_{G_{\Q_q}},b|_{G_{\Q_q}}$.  Write its third coordinate
afterwards as $h_q$.  By the same defect calculation,
\[
              \delta_q=[h_q-h|_{G_{\Q_q}}]
                    \in H^1(G_{\Q_q},\widehat\Z(2)).
 \tag{A.40.16}
\]
The uniqueness of the abelian conjugating element and the
remaining central conjugation show that this class is independent
of the choices.  Replacing $s$ by $s+k\eta$ changes
$\delta_q$ by $-\operatorname{res}_q(k\eta)$.
This proves the exact equality (A.40.5).  The isomorphism
(A.40.4) chooses a unique $k$ which kills the two indicated
comparison coordinates.  Its resulting global section class
is independent of the initially chosen lift and gives the
distinguished lift in the statement.
\end{proof}

The image in (A.40.3) is strictly smaller than the full space
of local central classes.  For an explicit coordinate proving
this, (A.40.1) and (A.40.13) give
\[
 H^1(G_\Q,\Z_5(2))=0,
 \qquad H^1(G_{\Q_{11}},\Z_5(2))
             =\Z_5/120\Z_5\simeq\Z/5.
 \tag{A.40.17}
\]
A local family supported in this nonzero coordinate cannot be
the restriction of a global central twist.  This calculation
does not assert that such an arbitrary family arises as
(A.40.16) from an ordinary unit-equation pair.  The unresolved
assertion is precisely whether every pair in $\mathcal G$
has its actual family (A.40.16) in the explicit twenty-four
element image (A.40.3).  A.38 alone gives no such membership:
its primitive solves the global cup equation, whereas (A.40.16)
also prescribes its local central cohomology classes.
